\documentclass[10pt,a4paper]{amsart}
\usepackage[margin=19mm]{geometry}
\usepackage{amsmath,amssymb,mathtools}
\usepackage[scr=rsfs,cal=euler]{mathalfa}
\usepackage{xfrac}
\usepackage[unicode,hidelinks,bookmarks=false]{hyperref}
\newcommand{\ie}{i.e.\ }
\newcommand{\cf}{cf.\ }
\newcommand{\resp}{respectively\ }
\newcommand{\Subsection}{\subsection}
\providecommand{\nobreakdash}{\nobreak\hbox{-}}
\setlength{\emergencystretch}{2em}
\multlinegap0pt
\usepackage{bm}
\usepackage{bbm}
\usepackage{dsfont}
\usepackage{xspace}
\usepackage{cancel}
\usepackage{mathtools}
\usepackage{amsthm}
\makeatletter
\@ifundefined{theorem}{\newtheorem{theorem}{Theorem}}{}
\@ifundefined{corollary}{\newtheorem{corollary}[theorem]{Corollary}}{}
\@ifundefined{lemma}{\newtheorem{lemma}[theorem]{Lemma}}{}
\@ifundefined{proposition}{\newtheorem{proposition}[theorem]{Proposition}}{}
\makeatother
\title[Completeness of the Model Space Does Not Force Regularity fo]{Completeness of the Model Space Does Not Force Regularity for Infinite-Dimensional Lie Groups}
\author{Zongjian Han, Fungo Saluta}
\date{Source: arXiv:2607.28557v2 (CC BY 4.0)}
\begin{document}
\maketitle

\noindent\emph{Source and licence.} Completeness of the Model Space Does Not Force Regularity for Infinite-Dimensional Lie Groups. Version arXiv:2607.28557v2 (CC BY 4.0), distributed under the Creative Commons Attribution 4.0 International licence, as recorded in the arXiv licence metadata for this exact version and shown on the abstract page https://arxiv.org/abs/2607.28557v2. Full rights notice: licenses/arxiv-2607.28557-CC-BY-4.0.txt.\par\smallskip

\noindent\textit{Editorial scope.} Counted unit: the Corollary whose source label is \texttt{cor:evolution-domain} in the article (source0.tex lines 1009--1026, source label \texttt{cor:evolution-domain}), delivered together with the article's own complete original proof of it (source0.tex lines 1028--1086, one \texttt{proof} environment of 59 lines). No mathematics is written, repaired or re-derived: statement and proof are the article's own text line for line. Required context retained uncounted: prop:topological-algebra (lines 555--564); lem:smooth-pulses (lines 860--868); eq:pulse-derivative-estimate (lines 879--884); prop:local-evolution (lines 906--922); thm:no-evolution (lines 942--951); cor:all-k (lines 993--999). Every definition, display and label of those lines is retained unchanged. Omitted from this package and not redistributed: 1--554, 565--859, 869--878, 885--905, 923--941, 952--992, 1000--1008, 1027--1027, 1087--1474. Nothing inside a retained excerpt is omitted or altered: every retained body is delivered exactly as the article's lines are written, so no omission receipt of any kind accompanies this package. Citations: the retained excerpts carry no citation command at all, so no bibliography entry of the article is transcribed and no reference is redistributed. Reference adaptation: none was needed and none was made; every reference inside the delivered unit resolves inside it, so no unresolved reference or citation is produced. Printed slips kept literal: no printed word, symbol, hypothesis or number of the counted statement or of its proof was altered. Layout and class: the article's own class is \texttt{article} with packages including inputenc, geometry, amsmath, amssymb, amsthm, mathtools, microtype, enumitem, hyperref, cleveref; the wrapper is the lane's pinned \texttt{amsart} editorial wrapper at 10pt on A4, which restates the theorem-like environments the retained text needs and adds the article's own macro definitions that the retained text needs (none), with extra wrapper packages (bm, bbm, dsfont, xspace, cancel, mathtools). The delivered numbering is the unit's own. Assets: no retained span contains a figure environment, a graphics-inclusion command, a PDF-page inclusion, a table or a TikZ picture; no image file is copied into this package, the package declares no asset dependency and no image file is redistributed with it. Licence: version arXiv:2607.28557v2 of this article is distributed under the Creative Commons Attribution 4.0 International licence, as recorded in the arXiv licence metadata for this exact version and shown on the abstract page https://arxiv.org/abs/2607.28557v2. A full rights notice is delivered at licenses/arxiv-2607.28557-CC-BY-4.0.txt. Characters: every retained excerpt is pure ASCII, so the delivered engine, XeLaTeX with its default Latin Modern text font, reports no missing character.
\begin{proposition}
\label{prop:topological-algebra}
The space $(\mathcal B_A,\tau_{\mathrm f})$ is complete, and multiplication
\[
  \mathcal B_A\times\mathcal B_A\to\mathcal B_A,
  \qquad (x,y)\mapsto xy,
\]
is jointly continuous.  Every bounded subset of $\mathcal B_A$ lies in a
finite-dimensional subspace.
\end{proposition}

\begin{lemma}[Smooth pulse family]
\label{lem:smooth-pulses}
The curves $\gamma_\lambda$ belong to $C^\infty([0,1],W)$ and
\[
  \gamma_\lambda\longrightarrow0
  \quad\text{in }C^\infty([0,1],W)
  \quad(\lambda\downarrow0).
\]
\end{lemma}

\begin{equation}
  \sup_{t\in I_n}\|\gamma_\lambda^{(r)}(t)\|
  \leq C_r\lambda\varepsilon_n\delta_n^{-r-1}
  =C_r\lambda 2^{-n^3+n(r+1)}.
  \label{eq:pulse-derivative-estimate}
\end{equation}

\begin{proposition}[Unique evolution before the accumulation time]
\label{prop:local-evolution}
The formula
\begin{equation}
  U_\lambda(t):=
  h_{n,\lambda}\left(\frac{t-t_{n-1}}{\delta_n}\right)
  P_{n-1}(\lambda),
  \qquad t\in I_n,
  \label{eq:local-evolution}
\end{equation}
defines the unique $H_A$-valued $C^1$ solution on $[0,1)$ of
$U_\lambda'=\gamma_\lambda U_\lambda$ with $U_\lambda(0)=1$.  This
solution is smooth and satisfies
\[
  U_\lambda(t_n)=P_n(\lambda).
\]
\end{proposition}

\begin{theorem}[Absence of a global $C^1$ evolution]
\label{thm:no-evolution}
For every $0<\lambda\leq1$, the initial-value problem
\[
  U'(t)=\gamma_\lambda(t)U(t),
  \qquad U(0)=1,
\]
has no $C^1$ solution with values in $G_A$, and hence none with values in
$H_A$.
\end{theorem}

\begin{corollary}
\label{cor:all-k}
For every $k\in\mathbb N\cup\{\infty\}$, the groups $H_A$ and $G_A$ are
not $C^k$-semiregular.  Every neighbourhood of the zero curve in
$C^k([0,1],W)$ contains a smooth curve admitting neither an $H_A$-valued nor
a $G_A$-valued right evolution.
\end{corollary}

\begin{corollary}[Instability of the evolution domain]
\label{cor:evolution-domain}
Let
\[
  \mathcal D_W:=\{\gamma\in C^\infty([0,1],W):
  U'=\gamma U,\ U(0)=1,
  \text{ has an $H_A$-valued $C^1$ solution}\}.
\]
Then $0\in\mathcal D_W$, the zero curve is not an interior point of
$\mathcal D_W$, and $\mathcal D_W$ is not closed in
$C^\infty([0,1],W)$.  The endpoint map, with $\mathcal D_W$ carrying the
subspace topology,
\[
  \operatorname{evol}_W:\mathcal D_W\to H_A,
  \qquad \gamma\mapsto U_\gamma(1),
\]
is well defined and is not continuous at $0$.
\end{corollary}

\begin{proof}
The uniqueness calculation in the proof of
Proposition~\ref{prop:local-evolution} applies to any two $H_A$-valued
$C^1$ solutions of the same right-evolution equation; hence
$\operatorname{evol}_W$ is well defined.  Moreover, a $C^1$ solution for a
smooth control is automatically smooth.  Indeed, in the affine chart
$H_A\cong A$, write $U=1+u$.  The equation becomes
\[
  u'=\gamma+\gamma u.
\]
Since multiplication is continuous bilinear, its right-hand side is $C^m$
whenever $u$ is $C^m$; induction gives $u\in C^m$ for every $m$.
Thus $\mathcal D_W$ is precisely the usual smooth right-evolution domain
restricted to $C^\infty([0,1],W)$.

The inclusion $0\in\mathcal D_W$ is immediate, and
Corollary~\ref{cor:all-k} shows that $0$ is not an interior point.  Fix
$\lambda>0$.  For $N\geq1$, define
\[
  \gamma_{\lambda,N}(t):=
  \begin{cases}
    \gamma_\lambda(t),&0\leq t\leq t_N,\\
    0,&t_N\leq t\leq1.
  \end{cases}
\]
The pulse construction is identically zero near every joining point, so
$\gamma_{\lambda,N}$ is smooth.  It admits the global evolution
\[
  U_{\lambda,N}(t):=
  \begin{cases}
    U_\lambda(t),&0\leq t\leq t_N,\\
    P_N(\lambda),&t_N\leq t\leq1.
  \end{cases}
\]
For every derivative order $r$, estimate
\eqref{eq:pulse-derivative-estimate} shows that the $C^r$-norm of the omitted
tail tends to zero.  Hence
$\gamma_{\lambda,N}\to\gamma_\lambda$ in
$C^\infty([0,1],W)$.  The limit does not belong to $\mathcal D_W$ by
Theorem~\ref{thm:no-evolution}; thus $\mathcal D_W$ is not closed.

For endpoint continuity, take $\lambda_N:=2^{-N}$ and
$\xi_N:=\gamma_{\lambda_N,N}$.  The uniform bounds in
Lemma~\ref{lem:smooth-pulses} imply $\xi_N\to0$ in
$C^\infty([0,1],W)$, while
\[
  \operatorname{evol}_W(\xi_N)=P_N(\lambda_N).
\]
The set $\{P_N(\lambda_N):N\geq1\}$ is not bounded: if it were bounded,
Proposition~\ref{prop:topological-algebra} would place it in a fixed
finite-dimensional subspace, whereas
\[
  \pi_N(P_N(\lambda_N)-1)
  =\lambda_N^N\varepsilon_1\cdots\varepsilon_Nw_N\neq0
\]
introduces a new homogeneous degree for every $N$.  Therefore the endpoints
do not converge to the identity, and $\operatorname{evol}_W$ is not
continuous at $0$.
\end{proof}

\end{document}
